\section*{Цель работы}
Ознакомление с методикой поверки аналоговых электромеханических измерительных приборов и определение их основных метрологических характеристик.
 
\section*{Задание}
\begin{enumerate}
    \item Ознакомиться с имеющейся на рабочем месте аппаратурой и получить от преподавателя конкретное задание для выполнения лабораторной работы.
    \item Определить основную погрешность комбинированного измерительного прибора (тестера) в следующих режимах работы:
    \begin{itemize}
        \item вольтметра постоянного тока;
        \item вольтметра переменного тока;
        \item миллиамперметра постоянного тока.
    \end{itemize}
    \item Для каждого из поверяемых режимов построить на одном графике зависимости относительной и приведенной погрешностей прибора от показаний поверяемого прибора.
    \item Определить амплитудно-частотную характеристику (АЧХ) вольтметра переменного тока. Построить график АЧХ, определить рабочую полосу частот вольтметра.
    \item На основе анализа полученных данных сделать вывод о соответствии поверяемого прибора его классу точности.
\end{enumerate}

\section*{Спецификация применяемых средств измерений}
\begin{table}[H]
    \centering
    \caption{Перечень используемого оборудования}
    \label{tbl:equipment}
    \begin{tblr}{
        colspec={X[l] c c c},
        hlines, vlines,
    }
        Наименование & Тип & Класс точности & Пределы измерений \\
        \SetCell[r=2]{} Поверяемый прибор (тестер) & & & \\
        & & & \\
        \SetCell[r=2]{} Образцовый прибор & & & \\
        & & & \\
        Источник постоянного тока & & & \\
        Генератор сигналов НЧ & & & \\
    \end{tblr}
\end{table}

\section*{Схемы экспериментов}
\begin{figure}[H]
    \centering
    \includegraphics[width=0.7\textwidth]{figs/scheme_voltmeter.pdf}
    \caption{Схема поверки вольтметра}
    \label{fig:scheme_voltmeter}
\end{figure}

\begin{figure}[H]
    \centering
    \includegraphics[width=0.5\textwidth]{figs/scheme_ammeter.pdf}
    \caption{Схема поверки амперметра}
    \label{fig:scheme_ammeter}
\end{figure}

\section*{Порядок выполнения работы}
\subsection*{Поверка измерительного прибора (Задания 2, 3)}
\begin{enumerate}
    \item Для поверки вольтметра постоянного и переменного тока собрать схему, показанную на \cref{fig:scheme_voltmeter}.
    \item Для поверки миллиамперметра постоянного тока собрать схему, показанную на \cref{fig:scheme_ammeter}.
    \item Поверка проводится методом сличения для 5--10 точек, равномерно распределенных по диапазону измерений.
    \item Для каждой поверяемой отметки $x$ на шкале прибора необходимо снять показания образцового прибора $x_{о.ув}$ при плавном увеличении измеряемой величины до отметки $x$.
    \item Затем, не меняя положения указателя поверяемого прибора, снять показания образцового прибора $x_{о.ум}$ при плавном уменьшении измеряемой величины до отметки $x$.
    \item Результаты измерений для вольтметра постоянного тока занести в \cref{tbl:protocol_vdc}.
    \item Результаты измерений для вольтметра переменного тока занести в \cref{tbl:protocol_vac}.
    \item Результаты измерений для миллиамперметра постоянного тока занести в \cref{tbl:protocol_idc}.
\end{enumerate}

\subsection*{Определение АЧХ вольтметра переменного тока (Задание 4)}
\begin{enumerate}
    \item В схеме, показанной на \cref{fig:scheme_voltmeter}, в качестве источника питания $U$ используется низкочастотный генератор сигналов.
    \item Установить частоту генератора $f = 50$ Гц. Регулируя выходное напряжение генератора, установить указатель поверяемого вольтметра на оцифрованную отметку в диапазоне (0,5\dots0,8) от выбранного предела измерений. Зафиксировать это значение напряжения $U(50)$.
    \item Не изменяя выходного напряжения генератора, уменьшать его частоту до прекращения колебаний указателя. Зафиксировать минимальное значение частоты $f_1$.
    \item Плавно увеличивая частоту, снять 2--3 точки для построения АЧХ до частоты $f = 50$ Гц.
    \item Продолжая увеличивать частоту входного напряжения до 10\dots20 кГц, снять не менее 5--7 точек для построения АЧХ. Рекомендуется снимать показания чаще в тех областях, где изменения показаний вольтметра наиболее заметны.
    \item Результаты измерений занести в \cref{tbl:protocol_afc}.
\end{enumerate}